% analysis/pipeline-DE.tex
\section{AIM-Verarbeitungspipeline}
\subsection*{Motivation}
Alignment-Modellierung umfasst mehrere Gebiete:
\begin{itemize}
\item globale Koordinatensysteme (CRS),
\item geometrische Alignment-Darstellung,
\item physische Realisierung,
\item Messdaten und
\item Optimierung.
\end{itemize}
Ihre getrennte Behandlung führt häufig zu fragmentierten Modellen und
inkonsistenten Ergebnissen. Im AIM-Rahmen werden sie in einer einzigen
Verarbeitungspipeline vereinigt.

\subsection{Operatorkette}
\begin{definition}[AIM-Pipeline]
Der vollständige Modellierungs- und Optimierungsprozess wird als
Operatorkomposition beschrieben:
\[
\boxed{\mathcal P=\mathcal W^{-1}\circ\mathcal R\circ\mathcal M\circ\mathcal W\circ\mathcal C}
\]
\end{definition}
Die Operatoren sind:
\begin{itemize}
\item \(\mathcal C\): CRS-Transformation,
\item \(\mathcal W\): World--Track-Transformation,
\item \(\mathcal M\): Alignment-Modell (parametrisch oder hybrid),
\item \(\mathcal R\): Realisierungsoperator,
\item \(\mathcal W^{-1}\): Track--World-Transformation.
\end{itemize}

\subsection{Detaillierte Struktur}
Die Pipeline bildet Rohdaten auf realisierte Geometrie ab:
\[
\text{CRS-Daten}\xrightarrow{\mathcal C}\text{World-Koordinaten}
\xrightarrow{\mathcal W}(s,d)\xrightarrow{\mathcal M}\kappa(s)
\xrightarrow{\mathcal R}\gamma_{\mathrm{real}}(s)
\xrightarrow{\mathcal W^{-1}}\text{World-Geometrie}.
\]

\subsection{Integration des hybriden Modells}
Der Alignment-Modelloperator ist hybrid:
\[
\mathcal M:(\mathbf T,\delta\kappa)\mapsto\kappa(s).
\]
Die Pipeline arbeitet somit zugleich mit strukturierten Parametern und lokalen
Korrekturen.

\subsection{Inverses Problem}
\begin{definition}[Pipeline-Inversion]
Für beobachtete Daten \(x_i\) lautet das inverse Problem:
\[
\min_{\mathbf x}\sum_i\|\mathcal P(\mathbf x)(s_i)-x_i\|^2.
\]
\end{definition}
Dies entspricht der Anpassung eines Modells, sodass die realisierte Geometrie
mit Beobachtungen übereinstimmt.

\subsection{Einbettung der Optimierung}
Die Pipeline ist in ein Optimierungsproblem eingebettet:
\[
\min_{\mathbf x}J\bigl(\mathcal P(\mathbf x)\bigr).
\]
Die Zielfunktion wird an der realisierten Geometrie und nicht am parametrischen
Modell bewertet.

\subsection{Jacobi-Struktur}
Die Ableitung der Pipeline folgt aus der Kettenregel:
\[
\frac{\partial\mathcal P}{\partial\mathbf x}
=D\mathcal W^{-1}\cdot D\mathcal R\cdot D\mathcal M\cdot D\mathcal W.
\]
Die Komponenten besitzen spezifische Eigenschaften:
\begin{itemize}
\item \(D\mathcal M\): dünn besetzt und strukturiert,
\item \(D\mathcal R\): Glättungsoperator,
\item \(D\mathcal W\): nichtlineare Projektion,
\item \(D\mathcal C\): lineare Transformation.
\end{itemize}

\subsection{Rechnerische Folgen}
Die Pipeline weist dünne Besetzung auf Modellebene, Nichtlinearität auf
Projektionsebene, Glättung auf Realisierungsebene und Skalierungseffekte auf
CRS-Ebene auf. Effiziente Implementierung erfordert hybride Darstellungen,
blockstrukturierte Löser und LOD-abhängige Bewertung.

\subsection{Interpretation nach Detaillierungsgrad}
Auf verschiedenen Maßstäben dominieren verschiedene Pipeline-Anteile:
\begin{itemize}
\item großer Maßstab: CRS und World-Geometrie,
\item mittlerer Maßstab: parametrisches Alignment,
\item kleiner Maßstab: Korrekturen und Realisierung.
\end{itemize}

\subsection{Zentrale Erkenntnis}
\begin{proposition}
Alignment-Modellierung ist keine einzelne Abbildung, sondern eine Komposition
geometrischer, physischer und rechnerischer Operatoren.
\end{proposition}

\subsection{Verbindung zu AIM}
\begin{summary}
Der AIM-Rahmen wird grundlegend durch seine Operatorpipeline bestimmt. Sie
integriert Koordinatensysteme, Geometrie, Physik und Optimierung in ein
kohärentes Modell und ermöglicht konsistente Datenintegration, physisch
bedeutsame Optimierung und skalierbare Berechnung.

AIM überführt Alignment-Modellierung damit von einer Sammlung isolierter
Verfahren in ein strukturiertes rechnerisches System.
\end{summary}
